\documentclass[11pt]{article}
\usepackage[margin=0.8in]{geometry}
\usepackage{array}
\usepackage{longtable}
\usepackage{verbatim}
\usepackage[T1]{fontenc}
\usepackage{lmodern}

\title{RISCV\_CNN Final Project Report}
\author{Student ID / Name: \texttt{studID\_name}}
\date{May 6, 2026}

\begin{document}
\maketitle

\section{Degree of Completion}

This implementation completes the required standalone RISCV\_CNN final project design in a single Verilog source file. The top-level module is:

\begin{verbatim}
RISCV_CNN(clk, rstn, tc, mode, st, seven_seg, anode)
\end{verbatim}

The completed functions are:
\begin{itemize}
  \item RV32 in-order pipelined CPU subset: \texttt{lw}, \texttt{sw}, \texttt{add}, \texttt{sub}, \texttt{addi}, \texttt{beq}, and signed \texttt{blt}.
  \item Byte-addressed PC and CPU data accesses, with instruction index \texttt{pc[31:2]} and data word index \texttt{addr[11:2]}.
  \item CPU host program that polls the start/config word, clears the start bit, computes Bias0/Bias1, stores them to shared memory, and launches the CNN coprocessor.
  \item Two-stage 3x3 CNN coprocessor using signed Q1.6 feature values and weights.
  \item Shared 1024-word data memory with CPU/tester access on Port A and CNN access on Port B.
  \item Completion protocol: after the CNN finishes, \texttt{DataMemory[13] = (512 << 1) | 1 = 1025}.
  \item Local tester/display wrapper for standalone board and simulation verification.
\end{itemize}

\section{Architecture}

\begin{verbatim}
        tc/mode/st/rstn/clk
              |
              v
 +-----------------------------+
 |         RISCV_CNN top       |
 |                             |
 | +----------+   Port A       |
 | |  Local   |----mux----+    |
 | | Tester   |           |    |
 | +----------+           v    |
 |                  +----------+
 | +----------+     | Shared   |     Port B     +----------+
 | | RV32 CPU |<--->| 1024x32  |<-------------->|   CNN    |
 | +----------+     | DataMem  |                | Coproc   |
 |                  +----------+                +----------+
 |                             |
 |                    seven-segment display
 +-----------------------------+
\end{verbatim}

The compute core is separated from the local tester/display logic so that the TA-provided test circuit can be integrated later by replacing the wrapper-side behavior while keeping the CPU, shared memory, and CNN coprocessor unchanged.

\section{Memory Map}

All PDF memory addresses are treated as word addresses internally. CPU load/store instructions remain standard RV32 byte-addressed accesses.

\begin{longtable}{|>{\raggedright\arraybackslash}p{1.4in}|>{\raggedright\arraybackslash}p{4.7in}|}
\hline
\textbf{Word Address} & \textbf{Use} \\
\hline
0--2 & Packed Weight0 3x3 kernel bytes \\
\hline
3--5 & Packed Weight1 3x3 kernel bytes \\
\hline
6--7 & Student command/status area; word 6 launches CNN \\
\hline
12--13 & Configuration/start word and final status word \\
\hline
14--15 & Bias0 and Bias1 generated by the CPU host program \\
\hline
16--271 & Packed input feature map \\
\hline
272--511 & First convolution intermediate feature map \\
\hline
512--959 & Final output feature map \\
\hline
960--1023 & Reserved; design avoids writes to this region \\
\hline
\end{longtable}

Feature maps are packed row-major with four signed 8-bit Q1.6 values in each 32-bit word. The lane order is big-endian: \texttt{[31:24]}, \texttt{[23:16]}, \texttt{[15:8]}, then \texttt{[7:0]}.

\section{CPU Design}

The CPU is an in-order pipelined RV32 subset processor with a small instruction ROM. It uses IF/ID, ID/EX, EX/MEM, and MEM/WB pipeline registers, conservative RAW hazard stalls, and branch flushing when \texttt{beq} or signed \texttt{blt} resolves in the EX stage. Register \texttt{x0} is hardwired to zero. The PC is byte-addressed and advances by 4 for normal instructions. Branch immediates are sign-extended, and \texttt{blt} uses signed comparison.

The ROM program performs the required host-control sequence:
\begin{enumerate}
  \item Poll \texttt{DataMemory[12]} until nonzero.
  \item Clear the start bit by subtracting 1 and storing the updated word back to \texttt{DataMemory[12]}.
  \item Execute the Figure 7 style bias-generation sequence using corrected RV32 byte addresses. The input base word 16 is accessed by byte address 64.
  \item Store Bias0 to word 14 and Bias1 to word 15.
  \item Write command word 6 to launch the CNN coprocessor.
\end{enumerate}

\section{CNN Coprocessor Design}

The CNN coprocessor reads the feature-map size from \texttt{DataMemory[12][6:1]}. It loads both 3x3 kernels from packed weight words and performs two convolution stages:

\begin{itemize}
  \item Stage 1 reads the input feature map at word 16 and writes the \((N-2)\times(N-2)\) intermediate map at word 272.
  \item Stage 2 reads the intermediate map and writes the \((N-4)\times(N-4)\) output map at word 512.
\end{itemize}

The coprocessor uses a row-MAC datapath with three signed 8x8 multipliers per kernel row. If a 3-value window crosses a packed 32-bit word boundary, the datapath fetches the next word and selects the needed lane bytes. After each 3x3 convolution output, the selected bias is added, the Q2.12 accumulator is rounded to Q1.6 with round-to-nearest ties-to-even, and the result is saturated to the signed 8-bit range \([-128, 127]\).

\section{Verification}

Behavioral simulation was run with the local self-checking testbench \texttt{tb\_riscv\_cnn.v}. The testbench checks CPU instruction behavior, CNN outputs against a reference model, fixed-point edge cases, and memory safety assertions.

\begin{longtable}{|>{\raggedright\arraybackslash}p{1.1in}|>{\raggedright\arraybackslash}p{1.1in}|>{\raggedright\arraybackslash}p{1.6in}|>{\raggedright\arraybackslash}p{2.0in}|}
\hline
\textbf{Testcase} & \textbf{N} & \textbf{Cycles} & \textbf{Result} \\
\hline
TC0 & 10 & 2080 & PASS \\
\hline
TC1 & 11 & 2727 & PASS \\
\hline
TC3 & 16 & 6892 & PASS \\
\hline
TC7 & 31 & 31887 & PASS \\
\hline
TC8 & 32 & 34012 & PASS \\
\hline
\end{longtable}

The final RTL simulation transcript ended with:

\begin{verbatim}
ALL TESTCASES PASSED
\end{verbatim}

Evidence files generated by the verification flow:
\begin{itemize}
  \item RTL simulation transcript: \texttt{xsim.log}
  \item RTL waveform database: \texttt{tb\_riscv\_cnn\_sim.wdb}
  \item Vivado build log: \texttt{vivado\_build.log}
  \item Timing report: \texttt{reports/timing\_impl.rpt}
  \item Utilization report: \texttt{reports/utilization\_impl.rpt}
  \item Route status report: \texttt{reports/route\_status.rpt}
\end{itemize}

\section{Synthesis and Implementation Result}

Vivado implementation completed successfully for device \texttt{xc7a35tcsg324-1} with a 10 ns system clock constraint.

\begin{longtable}{|>{\raggedright\arraybackslash}p{2.0in}|>{\raggedright\arraybackslash}p{2.0in}|>{\raggedright\arraybackslash}p{1.4in}|}
\hline
\textbf{Metric} & \textbf{Result} & \textbf{Utilization} \\
\hline
Slice LUTs & 1513 / 20800 & 7.27\% \\
\hline
Slice Registers & 1613 / 41600 & 3.88\% \\
\hline
Block RAM Tile & 1 / 50 & 2.00\% \\
\hline
DSPs & 1 / 90 & 1.11\% \\
\hline
Setup timing & 0 failing endpoints, WNS 0.094 ns & PASS \\
\hline
Hold timing & 0 failing endpoints, worst slack 0.082 ns & PASS \\
\hline
Pulse width timing & 0 failing endpoints, worst slack 4.500 ns & PASS \\
\hline
Routing & 0 routing errors & PASS \\
\hline
Bitstream & \texttt{FP\_studID.bit} generated & PASS \\
\hline
\end{longtable}

The final bitstream DRC reported 0 errors and 3 warnings. The warnings are DSP pipelining recommendations and do not prevent bitstream generation or timing closure.

\section{Bonus Design}

For the pipeline CPU bonus, the top-level system instantiates \texttt{rv32\_pipelined\_cpu}. The previous multi-cycle CPU source is kept in the same file for reference, but it is not instantiated by \texttt{RISCV\_CNN}. The pipeline uses conservative hazard handling, so dependent instructions are stalled until the producing instruction has safely written back. This keeps the CPU control simple and robust while satisfying the pipelined CPU design requirement.

\section{Constraints Note}

The project PDF lists reset as pin S6, but Vivado reports S6 as an illegal package pin for \texttt{xc7a35tcsg324-1}. The included constraint file therefore maps \texttt{rstn} to P15, matching previous EGO1 lab constraint usage. The remaining top-level pins follow the PDF pin-facing interface.

\section{Conclusion}

The design completes all required CPU instructions, CPU bias generation, memory-mapped CNN launch, two-stage CNN computation, fixed-point rounding/saturation, and final output-address/status protocol. Behavioral simulation passed all local reference tests, implementation timing passed at the 10 ns clock constraint, and the bitstream was generated successfully.

\end{document}
